\documentclass[12pt]{article}
\usepackage[utf8]{inputenc}
\usepackage[T1]{fontenc}
\usepackage[margin=1in]{geometry}
\usepackage{setspace}
\usepackage{xcolor}
\usepackage[colorlinks=true, urlcolor=blue!60!black, linkcolor=blue!60!black]{hyperref}
\setlength{\parindent}{0em}
\setlength{\parskip}{0.45em}
\renewcommand{\baselinestretch}{1.0}
\begin{document}
\begin{center}
{\Large \textit{Rucho v. Common Cause}, 588 U.S. 684 (2019)\par}
\vspace{0.3em}
{\normalsize Case Brief --- One-Page Edition\par}
\end{center}

\textbf{Memory Jogger.} Partisan gerrymandering claims are nonjusticiable in federal court: no neutral standard for how much is too much (\href{https://www.courtlistener.com/opinion/4633469/rucho-v-common-cause/}{CourtListener}).

\textbf{Facts.} North Carolina Republicans adopted ``Partisan Advantage'' as a criterion for a 10--3 congressional map; Maryland Democrats moved some 700,000 voters to flip the Sixth District and win 7 of 8 seats.

\textbf{Procedural History.} Three-judge courts struck down North Carolina's map under equal protection, the First Amendment, and Article I, and Maryland's under the First Amendment; the States appealed directly.

\textbf{Votes.} Vacated 5--4. Roberts wrote for the Court; Kagan, with Ginsburg, Breyer, and Sotomayor, dissented.

\textbf{Holding.} Partisan gerrymandering claims under any provision present political questions with no judicially manageable standards; federal courts must dismiss them.

\textbf{Reasoning.} Gerrymandering is as old as the Republic, and the Elections Clause assigns oversight to Congress. Because some partisanship is lawful, the question is one of degree, and the claims rest on proportional representation, which the Constitution does not require; one-person, one-vote and racial claims remain reviewable. Kagan answered that lower courts had a workable test using simulated maps---``This much is too much''---and that politicians will not reform themselves.

\textbf{Examples.} \emph{Same-Side:} a state legislative map locking in a supermajority. \emph{Opposite-Side:} state-court enforcement of state constitutions. \emph{Fence-Sitter:} race and party intertwined.

\textbf{Critique.} Conservatives praised deference to the Elections Clause, and John McGinnis called it ``silent originalism''; election-law scholars call it abdication that greenlights technological entrenchment.

\textbf{Key Quotation.} ``Excessive partisanship in districting leads to results that reasonably seem unjust. But the fact that such gerrymandering is `incompatible with democratic principles,' \textit{Arizona State Legislature}, 576 U.\,S., at 791 (brackets omitted), does not mean that the solution lies with the Federal Judiciary. We conclude that partisan gerrymandering claims present political questions beyond the reach of the federal courts.'' --- Chief Justice Roberts, 588 U.S. at 718.

\end{document}
